\begin{figure}[ht]
\centering
\begin{tikzpicture}[x=1.0cm, y=1.0cm, font=\small, >=Stealth]

% Stage labels along the top
\node[font=\scriptsize] at (1.5, 3.4) {original $\mat{A}$};
\node[font=\scriptsize] at (5.5, 3.4) {after row swap $\mat{P}\mat{A}$};
\node[font=\scriptsize] at (9.5, 3.4) {eliminated $\mat{U}$};

% Stage 1: original, zero (small) pivot flagged
\node at (1.5, 2.0) {$\begin{pmatrix} 0 & 2 & 1 \\ 1 & 1 & 0 \\ 2 & 1 & 3 \end{pmatrix}$};
\node[engred, font=\scriptsize, anchor=north] at (1.5, 0.8) {pivot $=0$: swap needed};

% arrow
\draw[->, thick, enggray] (2.7, 2.0) -- (3.9, 2.0);

% Stage 2: after swapping row 1 with row 3 (largest |entry|)
\node at (5.5, 2.0) {$\begin{pmatrix} 2 & 1 & 3 \\ 1 & 1 & 0 \\ 0 & 2 & 1 \end{pmatrix}$};
\node[engblue, font=\scriptsize, anchor=north] at (5.5, 0.8) {largest entry to pivot};

% arrow
\draw[->, thick, enggray] (7.1, 2.0) -- (8.3, 2.0);

% Stage 3: upper triangular
\node at (9.5, 2.0) {$\begin{pmatrix} 2 & 1 & 3 \\ 0 & 0.5 & -1.5 \\ 0 & 0 & 7 \end{pmatrix}$};
\node[engamber, font=\scriptsize, anchor=north] at (9.5, 0.8) {$\mat{P}\mat{A} = \mat{L}\mat{U}$};

\node[anchor=north, font=\scriptsize, text width=12cm, align=center]
  at (5.5, 0.1) {
  A zero (or small) pivot forces a row exchange before elimination can
  proceed. The permutation matrix $\mat{P}$ records every swap, so the
  factorisation the algorithm actually produces is of the permuted
  matrix, $\mat{P}\mat{A} = \mat{L}\mat{U}$.
};

\end{tikzpicture}
\caption[Partial pivoting produces the permuted factorisation]{Partial
pivoting swaps the row with the largest available entry into the pivot
position before each elimination step, both to avoid a zero pivot and
to control the growth of round-off. The swaps are collected into a
permutation matrix $\mat{P}$, and the factorisation the algorithm
delivers is of the permuted matrix, $\mat{P}\mat{A} = \mat{L}\mat{U}$.
Applying the same permutation to the right-hand side, then solving
$\mat{L}\vect{y} = \mat{P}\vect{b}$ and $\mat{U}\vect{x} = \vect{y}$,
returns the solution of the original system.}
\label{fig:vol02:ch09:pivot-permutation}
\end{figure}
